\chapter{Propuesta científica y correspondencia con el proyecto de Ciencia de Frontera}
\label{ch:secihti_mapping}

\section{Problema, hipótesis y alcance de la propuesta actualizada}

La propuesta aborda la dificultad de recuperar movimiento y fisiología mediante
infraestructura inalámbrica cuando la señal útil se superpone con multitrayectoria,
deriva instrumental y cambios del entorno. Los avances de dc01 muestran que mejorar
potencia y retardos no garantiza predecir fase bruta, y que una corrección espectral
parcialmente predecible sí aporta concordancia bajo las condiciones documentadas.
El siguiente problema científico consiste en determinar qué transformaciones pueden
compensarse o hacerse invariantes sin perder la respuesta física de interés.

Se propone un gemelo digital con geometría métrica, modelo de observación y dos ramas
complementarias: una robusta para estructura relativa y otra coherente para dinámica.
La hipótesis central y H1--H6 se formalizan en la \cref{sec:updated_hypotheses}.
El contraste requiere intervenir físicamente el entorno y caracterizar la referencia
de adquisición; no se resuelve únicamente añadiendo capacidad neuronal al ajuste
estático. La aportación esperada es un procedimiento reproducible que relacione
calibración, representación, sensibilidad y transferencia con incertidumbre explícita.

El programa conserva Wi-Fi y LoRa como radios de interés. SDR, radar, visión y
sensores de referencia se emplearán para validar y supervisar, declarando qué modalidades
son necesarias durante inferencia. La aplicación humana se desarrollará después del
reflector y el fantoma. El trabajo de laboratorio y la demostración comunitaria tendrán
alcances propios; no se atribuye validación clínica ni despliegue formal a la Etapa 1.

\section{Correspondencia científica y administrativa}

El programa aprobado se organiza en diseño, implementación experimental e integración
científica. La actualización precisa la secuencia técnica y los criterios de evidencia;
no constituye por sí misma una modificación del presupuesto o de compromisos autorizados.

\begin{table}[htbp]
\centering\small
\caption{Correspondencia entre etapas del proyecto y evidencia del programa actualizado.}
\label{tab:proposal_updated_mapping}
\begin{tabularx}{\textwidth}{P{3.0cm}L L}
\toprule
Etapa & Trabajo doctoral & Evidencia y entregables\\
\midrule
Diseño & L1/L2/L4; S4.7a, referencia, G0--G4 y Esc. E0--E1. & Especificación geométrica,
protocolo, comparaciones estáticas y caracterización de repetibilidad/sensibilidad.\\
Implementación & L3/L4/L5; Esc. E2--E3 y operadores multirradio. & Fantoma, referencia
sincronizada y estudio humano autorizado; datos con procedencia y particiones.\\
Integración & L5/L6; Esc. E4--E5, inversión y diseño de red. & Límites de separabilidad,
transferencia, incertidumbre y preparación del piloto según madurez demostrada.\\
\bottomrule
\end{tabularx}
\end{table}

Las etapas no implican que toda actividad de una fila deba concluir en los primeros
diez meses. El cronograma del \cref{ch:wbs_stage1} distingue preparación, prueba
preliminar y validación posterior. En particular, varios usuarios y transferencia
externa se condicionarán por la evidencia acumulada y los acuerdos necesarios.

\section{Productos científicos y criterios de revisión}

El primer producto será un banco comparativo que conserve resultados brutos, relativos
y diagnósticos oráculo separados, con particiones espaciales reproducibles. El segundo
será un protocolo metrológico que registre escena, antenas, relojes y desplazamiento
local, y evalúe G0--G4 sin confundir geometría con ajuste material. El tercero será
un conjunto preliminar de mediciones con reflector y fantoma, seguido de datos humanos
cuando se cumplan las condiciones técnicas y éticas. El cuarto será la comparación
de ramas y métodos de referencia bajo presupuestos de información explícitos.

Cada producto deberá contener la pregunta contrastada, configuración, artefactos de
reproducción, distribución de errores, incertidumbre y limitaciones. La disponibilidad
de una malla, un modelo entrenado o una demostración visual no basta para declarar
validado el producto. El criterio científico de revisión será que otro grupo pueda
reconstruir la correspondencia entre estado físico, señal adquirida, predicción y
resultado evaluado.

\section{Viabilidad y decisiones ante evidencia adversa}

La viabilidad se apoya en reutilizar el equipo reportado, priorizar referencia RF y
movimiento medido, y ejecutar S4.7a y el diseño del banco en paralelo. La inversión
en reconstrucción LiDAR se justificará por la comparación geométrica y la utilidad
multiescena; su adquisición no condiciona iniciar E0 ni una prueba local de E1 con
geometría simple verificada. La selección comercial se decidirá mediante especificación,
cotización y compatibilidad dentro de los rubros autorizados.

Una invariancia demasiado amplia se restringirá o se reservará a la rama de contexto.
Una referencia inestable motivará recalibración instrumental. Si el gemelo no reproduce
la sensibilidad mecánica, se investigarán geometría local, mecanismos de propagación
y observación antes de usarlo para inversión humana. La ausencia de separabilidad
motivará abstención o rediseño de enlaces. Estas decisiones proporcionan resultados
científicos útiles aun cuando una hipótesis de mejora no se confirme.

La colaboración prospectiva con UPNA y el Dr. Francisco Falcone se orientará a repetir
el protocolo canónico en otras bandas, sujeto a acuerdos, acceso y financiación.
La vinculación con Colmenas Guadalajara conserva su función de caracterización de sitio
y divulgación. Ninguna de estas actividades se presenta como ya ejecutada o como
sustituto de una validación experimental independiente.
